\documentclass[11pt]{article}
\usepackage[margin=1in]{geometry}
\usepackage{booktabs,amsmath,microtype,siunitx}
\usepackage[table]{xcolor}
\usepackage[colorlinks,linkcolor=blue!50!black,urlcolor=blue!50!black]{hyperref}
\usepackage{listings}
\lstset{basicstyle=\ttfamily\footnotesize,breaklines=true,frame=single,
        backgroundcolor=\color{gray!8},columns=fullflexible}

\title{Parallel optimization of the DisloCluster coupled march\\
       \large Measured bottlenecks, four fixes, two falsified hypotheses,
       and one claim retracted}
\author{DisloCluster development note}
\date{2 September 2026}

\begin{document}
\maketitle

\begin{abstract}
The coupled march spends its time in two solvers: a \emph{slow} step
integrating independent stiff ODEs at every mesh node (ZrMicro, OpenMP batch)
and a \emph{fast} step solving a steady reaction--diffusion FE problem
(MoDELib3 \texttt{DDomp}). Neither was using the machine, and in both cases the
reason was not the one the phase labels suggested. The slow step saturated at
$2.94\times$ on a 40-core host and ran \emph{slower than serial} at its
configured thread count; the fast step scaled only $1.68\times$, with its
largest phase mislabelled ``precond'' when it is in fact serial sparse-matrix
construction. Four changes follow from the measurements --- sharding the slow
batch across processes, removing a token round trip that consumed 47\,\% of the
slow step, caching the FE sparsity pattern across Newton iterations, and
correcting the fast step's thread count. Every one is verified
\textbf{bit-identical} against the unoptimized code. End to end on a
\SI{200}{\nano\meter} hexagonal march to \SI{1}{dpa} the wall time falls
\SI{130.4}{\second} $\rightarrow$ \SI{107.2}{\second}; the slow step alone
improves $2.71\times$ beyond sharding, and sharding was itself $1.79\times$.
One earlier claim in this note --- that FE assembly was serial --- was wrong and
is retracted in \S\ref{sec:retract}. The largest remaining cost is identified,
and it is not linear algebra.
\end{abstract}

\section{The machine, and why it matters}
All measurements are on a dual-socket Intel Xeon Gold 6230: \textbf{40 physical
cores, 80 logical threads}, Windows 10, with the 3-D solver running under WSL2.

One property of this host shapes every result below. Windows partitions 80
logical processors into \emph{two processor groups} of 40, and a process is
confined to one group unless it explicitly calls the group-affinity API:

\begin{lstlisting}
System.Environment.ProcessorCount  ->  40      (not 80)
process affinity mask              ->  1099511627775 = 2^40 - 1
\end{lstlisting}

So a single process can address at most half the machine, and asking for 80
threads oversubscribes 40 logical processors two-to-one.

\section{The slow step}

\subsection{Thread scaling: measured}
The slow step is embarrassingly parallel by construction --- freezing the mobile
species removes the only spatial coupling, leaving $N$ independent
19-dimensional systems with no neighbour state. Yet threads stopped paying
almost immediately. Fixed work of 20\,000 cases, three passes each, machine
idle:

\begin{center}
\begin{tabular}{rrrrr}
\toprule
processes & threads & total & cases/s & speedup\\
\midrule
1 & 1 & 1 & 1\,544 & $1.00\times$\\
1 & 8 & 8 & 4\,541 & $2.94\times$\\
1 & 16 & 16 & 4\,370 & $2.83\times$\\
1 & 40 & 40 & 2\,678 & $1.73\times$\\
1 & 80 & 80 & 1\,130 & $\mathbf{0.73\times}$\\
\midrule
2 & 8 & 16 & 8\,965 & $5.81\times$\\
5 & 8 & 40 & 14\,347 & $9.29\times$\\
10 & 4 & 40 & 16\,848 & $10.91\times$\\
\rowcolor{green!12}
20 & 4 & 80 & \textbf{20\,438} & $\mathbf{13.23\times}$\\
40 & 2 & 80 & 18\,469 & $11.96\times$\\
\bottomrule
\end{tabular}
\end{center}

At its configured 80 threads the batch ran \emph{slower than one thread}, and
the same 80 threads arranged as 20 processes of 4 are $18\times$ apart. The work
parallelizes; it just does not parallelize inside one process.

\subsection{Two hypotheses, both falsified}
Before changing anything, the two obvious mechanisms were built and measured.
Both are \textbf{ruled out}, and recording that is the point of this section:

\begin{center}
\begin{tabular}{lll}
\toprule
Hypothesis & Test & Result\\
\midrule
False sharing on \texttt{status[c]}, \texttt{outbuf[c]}
  & rebuild with \texttt{schedule(dynamic,64)}
  & $1.00$--$1.01\times$, no effect\\
SUNDIALS error-handler lock
  & remove all 20\,000 messages at source
  & scaling unchanged\\
\bottomrule
\end{tabular}
\end{center}

A third candidate --- the volume of progress text written to the driver's pipe
--- was also dismissed by measurement: \SI{0.69}{\mega\byte} per invocation at
\SI{200}{\nano\meter}, milliseconds.

The surviving inference comes from the data itself: \emph{processes scale and
threads do not}, which implicates a resource private to a process and shared by
its threads. Every case constructs a \texttt{std::ostringstream} and assigns a
\texttt{std::string} against one per-process heap, which fits, but this is
\textbf{not demonstrated} and is offered as a hypothesis only.

\subsection{Fix 1: sharding}
\emph{Batching} merges work to amortize per-call overhead; \emph{sharding}
splits it to escape a per-worker ceiling. The code already batched --- all
roughly 190\,000 cases into one process, precisely to amortize the
Windows-dominant process-spawn and DLL-load cost --- and that argument is
correct, but it does not say when to stop.
\texttt{run\_cpp\_solver\_batch} now splits the batch across concurrent
processes, each a self-contained batch file with its own \texttt{@BASE} header,
results scattered back by original index. Each shard is still a batch, so the
two ideas compose, and \texttt{\_SHARD\_MIN\_CASES} is where they balance:
below 2\,000 cases the work stays in one process, because splitting would
re-pay the overhead batching exists to avoid.

Three invariants are load-bearing: independence of the cases (guaranteed by the
frozen mobile field), reassembly by original index, and interrupt and failure
handling covering \emph{every} shard rather than the one being waited on.

\paragraph{Result} at 189\,533 nodes (157\,204 integrations after dedup):
\begin{center}
\begin{tabular}{lrrl}
\toprule
configuration & wall & s/node & identical\\
\midrule
1 process, 8 threads & \SI{52.3}{\second} & $2.76\times10^{-4}$ & ---\\
1 process, 16 threads & \SI{60.1}{\second} & $3.17\times10^{-4}$ & yes\\
\rowcolor{green!12}
auto shards, $20\times4$ & \textbf{\SI{29.2}{\second}} & $1.54\times10^{-4}$ & \textbf{yes}\\
\bottomrule
\end{tabular}
\end{center}

$1.79\times$ against the \emph{best} single-process configuration, bit-identical.

\subsection{Fix 2: the token round trip}
\label{sec:marshal}
With the solver ceiling removed, profiling the slow step showed the bottleneck
had moved out of C++ entirely --- 47\,\% of it was Python string handling:

\begin{center}
\begin{tabular}{lrrrr}
\toprule
phase & before & \% & after & \\
\midrule
dedup & \SI{0.61}{\second} & 1.5 & \SI{0.72}{\second} &\\
\texttt{build\_immobile\_cases} & \SI{5.31}{\second} & 13.2 & \textbf{\SI{0.00}{\second}} &\\
\texttt{\_batch\_lines} & \SI{13.64}{\second} & 33.9 & \textbf{\SI{2.53}{\second}} &\\
solver + parse & \SI{20.68}{\second} & 51.4 & \SI{11.46}{\second} &\\
\midrule
\rowcolor{green!12}
TOTAL & \SI{39.94}{\second} & & \textbf{\SI{14.72}{\second}} & $\mathbf{2.71\times}$\\
\bottomrule
\end{tabular}
\end{center}

Every case is the same $\sim$87-token material base with its 19 \texttt{y0}
slots overwritten. \texttt{build\_immobile\_cases} materialized all 87 per case;
\texttt{\_batch\_lines} then parsed all of them back into dictionaries and
diffed them across 157\,204 cases to rediscover which keys were shared --- the
split \texttt{build\_immobile\_cases} had known all along --- before emitting
the $\sim$19 that varied. Roughly 78\,\% of the formatted tokens existed only to
be discarded.

\texttt{build\_immobile\_cases} now returns an \texttt{ImmobileCases} holding
the base and the \texttt{y0} array and formatting only what varies; the batch
writer takes \texttt{base\_line()}/\texttt{delta\_lines()} directly. Indexing
still yields the full \texttt{--key=value} list, so every other consumer and the
single-case path are unchanged, and slicing returns an \texttt{ImmobileCases} so
the shard writer keeps working. Note the ``solver + parse'' row also halves,
because that path calls \texttt{\_batch\_lines} once per shard.

\paragraph{Key order is the protocol.} \texttt{collect\_solver\_args}
\emph{already} emits \texttt{y0\_0..y0\_18}, so a case dictionary overwrites
them \textbf{in place} and the token order is the base's, not base-then-\texttt{y0}.
A first attempt guarded against that case and fell back to the old path --- which
is to say it silently did nothing, and profiled identically to before. The order
is now reconstructed explicitly, and byte-identity is asserted at
$n = 1, 2, 500, 5000$, on an all-identical batch (exercising the empty-delta
guard) and on a sliced shard.

\section{The fast step}

\subsection{Phase split and thread scaling}
\texttt{DDomp} reports its own stage timers. The split is stable across a
ninefold change in problem size, which is what justifies optimizing against it:

\begin{center}
\begin{tabular}{lrr}
\toprule
phase & \SI{200}{\nano\meter} (21\,052 nodes) & \SI{1000}{\nano\meter} (189\,533 nodes)\\
\midrule
diffusion & 0.4\,\% & 0.3\,\%\\
assembly & 25.3\,\% & 25.6\,\%\\
\rowcolor{yellow!25}
``precond'' & \textbf{45.9\,\%} & \textbf{50.7\,\%}\\
krylov & 28.5\,\% & 23.4\,\%\\
\bottomrule
\end{tabular}
\end{center}

Whole-solver thread scaling is poor at both sizes and \emph{worst at the
default}:

\begin{center}
\begin{tabular}{lrrrrrrr}
\toprule
threads & 1 & 4 & 8 & 16 & 20 & 40 & 80\\
\midrule
\SI{200}{\nano\meter} wall [s] & 28.6 & 19.3 & 18.3 & 17.6 & 17.2 & \textbf{17.0} & 22.1\\
\SI{1000}{\nano\meter} wall [s] & --- & --- & --- & --- & \textbf{476.0} & 484.3 & 510.8\\
\bottomrule
\end{tabular}
\end{center}

$1.68\times$ from 1 to 40 threads at \SI{200}{\nano\meter}, and a 7\,\% spread
at \SI{1000}{\nano\meter}. Convergence was identical at every setting (9 Newton
iterations, same residual), so this is pure overhead.

\subsection{Per-phase scaling, and a retraction}
\label{sec:retract}
Resolving the scaling \emph{by phase} is what settles where the serial time
actually is (\SI{200}{\nano\meter}, with the pattern cache of
\S\ref{sec:pattern} in place):

\begin{center}
\begin{tabular}{lrrrr}
\toprule
threads & assembly & ``precond'' & krylov & total\\
\midrule
1 & 6.03 & 3.48 & 6.97 & 16.52\\
4 & 4.14 & 3.45 & 3.67 & 11.30\\
20 & 3.72 & 3.83 & 2.63 & 10.23\\
\midrule
scaling & $1.62\times$ & \cellcolor{red!15}$\mathbf{0.91\times}$ & $2.65\times$ & $1.61\times$\\
\bottomrule
\end{tabular}
\end{center}

\paragraph{Retraction.} An earlier revision of this note asserted that FE
assembly was serial and listed ``parallelize assembly'' as the second priority.
\textbf{That is wrong.} \texttt{BilinearWeakForm::globalTriplets} is already
parallelized, over explicit contiguous chunks concatenated in chunk order
specifically so the triplet sequence --- and hence the summed matrix --- stays
bit-identical to the serial one. It scales $1.62\times$ on 20 threads, which is
poor, but poor is not absent, and the recommendation to parallelize it was
advice to do something already done.

The phase that genuinely does not scale is \textbf{``precond''}, flat and
slightly negative at $0.91\times$, for the reason in \S\ref{sec:pattern}.
Krylov, expected to be the serial bottleneck, is in fact the \emph{best}-scaling
phase at $2.65\times$.

\subsection{Fix 3: cache the sparsity pattern}
\label{sec:pattern}
The dominant phase is \emph{not} preconditioning. The reaction solver is
constructed with \texttt{use\_directSolver = false}, so its preconditioner is
Eigen's diagonal one, which is $O(n)$ and free. The timer instead brackets
\texttt{computeFromTriplets}, whose real content is sparse-matrix
\emph{construction}: \texttt{A.setFromTriplets} over the full 84\,208-dof
matrix, a pass over its nonzeros to strike out the Dirichlet rows and columns,
then \texttt{A1.setFromTriplets}. Two sort-and-accumulate passes per Newton
iteration, four to nine iterations per solve, and \texttt{setFromTriplets} is
serial in Eigen.

Across the Newton loop that pattern is \emph{invariant} --- same mesh, same weak
forms, same element loops --- and only the values move. The surrounding class
already exploits exactly this reasoning for the Dirichlet selection matrix $T$
and its index map; the argument extends to the matrices themselves. The first
call now also records
\[
  \texttt{aSlot}[k]:\ \text{triplet } k \longmapsto \text{slot in } A,
  \qquad
  \texttt{a1FromA}[j]:\ \text{slot } j \text{ of } A \longmapsto \text{slot in } A_1,
\]
by binary search inside each compressed row, $O(\mathrm{nnz}\log\mathrm{nnz})$
\emph{once}; every later call zeroes the value arrays and accumulates in place,
skipping both sorts and every allocation. Values go in triplet order --- the
order \texttt{setFromTriplets} itself sums duplicates in --- so the refilled
matrix is bit-identical, and the cache is invalidated whenever the triplet count
or either dimension moves.

\begin{center}
\begin{tabular}{lrrr}
\toprule
\SI{200}{\nano\meter}, per fast solve & baseline & pattern cache & gain\\
\midrule
``precond'' & \SI{4.758}{\second} & \textbf{\SI{2.885}{\second}} & $1.65\times$\\
solve total & \SI{10.378}{\second} & \textbf{\SI{8.095}{\second}} & $1.28\times$\\
\bottomrule
\end{tabular}
\end{center}

Convergence is unchanged to every printed digit --- 4 iterations,
$c_{\mathrm{Error}} = 8.785110566792336\times10^{-7}$ before and after. The
residual cost is inherent: the first Newton iteration still builds the pattern,
so one full construction per solve remains, and that is why the phase still does
not scale.

\subsection{Fix 4: stop oversubscribing}
\texttt{qssa.py} launched \texttt{DDomp} through \texttt{wsl.exe} with no
environment at all, so it inherited 80 threads --- the worst setting measured at
both sizes. It now sets \texttt{OMP\_NUM\_THREADS=20}, overridable through
\texttt{DISLOCLUSTER\_DDOMP\_THREADS}, and forwards it via \texttt{WSLENV},
which WSL does not otherwise inherit.

\section{End-to-end result}
A \SI{200}{\nano\meter} hexagonal march to \SI{1}{dpa} (doses $0.01$, $0.1$,
$1.0$; 5 substeps per interval; \texttt{fem\_every} $=5$), run three times ---
baseline \texttt{DDomp} with the original Python, the pattern-cached
\texttt{DDomp}, and both plus the marshalling fix:

\begin{center}
\begin{tabular}{lrrr}
\toprule
 & fast solve (mean of 3) & slow step (sum) & march wall\\
\midrule
sharding only & \SI{22.9}{\second} & \SI{51}{\second} & \SI{130.4}{\second}\\
$+$ pattern cache & \SI{20.7}{\second} & \SI{47}{\second} & \SI{122.4}{\second}\\
\rowcolor{green!12}
$+$ marshalling & \SI{18.1}{\second} & \textbf{\SI{28}{\second}} & \textbf{\SI{107.2}{\second}}\\
\midrule
gain & $1.27\times$ & $\mathbf{1.82\times}$ & $\mathbf{1.22\times}$\\
\bottomrule
\end{tabular}
\end{center}

\texttt{march\_state.npz} is \textbf{bit-identical} across all three: max
absolute difference $0.000\times10^{0}$, over $4\times21\,052\times19$ values.

The end-to-end factor is smaller than the component factors, and it is worth
saying why rather than quoting the best number. At \SI{200}{\nano\meter} the
mobile solve is only $\sim$60\,\% of a \texttt{DDomp} invocation, and the case
is small enough that the slow step is not dominant. The per-component summary is
the honest one: slow step $2.71\times$ (on top of sharding's $1.79\times$),
fast-step linear algebra $1.28\times$, fast-step wall $1.27\times$.

\section{The largest remaining cost is not the linear algebra}
At \SI{1000}{\nano\meter} the stage timers total \SI{237}{\second} against a
\SI{476}{\second} wall. \textbf{Half the fast step lies outside the mobile solve
entirely} --- mesh reading, trial-function construction, FE data structures and
\texttt{evl} output. At \SI{200}{\nano\meter} the same gap is small, so it grows
faster than the solve it accompanies.

The coupled march invokes \texttt{DDomp} 64 times per \SI{40}{dpa} run and
rebuilds all of it from scratch every time: on the order of \SI{4}{\hour} of
pure re-initialization inside a \SI{15.4}{\hour} run. It needs no numerical work
--- only keeping the solver resident across fast solves, or caching the
assembled structures --- and it would additionally let the sparsity-pattern
cache of \S\ref{sec:pattern} amortize across all 64 solves instead of being
rebuilt inside each one.

\paragraph{Caveat.} \SI{239}{\second} is a \emph{residual} (wall minus timers),
not a breakdown. The diagnostic run that would apportion it was truncated by
129\,569 lines of mesh-progress output. Before budgeting four hours of saving,
that needs one clean run with the progress output silenced.

\section{Priorities}
\begin{enumerate}
  \item \textbf{Eliminate repeated setup} --- $\sim$50\,\% of fast-step wall at
        \SI{1000}{\nano\meter}, and not an OpenMP problem. Largest by a wide
        margin. Measure the breakdown first.
  \item \textbf{Assembly efficiency} --- already parallel, but only
        $1.62\times$ on 20 threads. The question is why (chunk imbalance,
        memory bandwidth, or per-element cost variance), not whether to
        parallelize it.
  \item \textbf{The first Newton iteration's matrix build} --- what remains of
        ``precond''. Subsumed by item 1 if the solver becomes resident.
  \item \textbf{\texttt{dedup\_rtol}} --- set to $10^{-8}$ against an integrator
        \texttt{rtol} of $10^{-6}$; dedup measures $\times1.21$. Loosening it
        cuts slow work proportionally with error bounded by a tolerance already
        accepted. Untested, and the cheapest thing on this list.
  \item \textbf{Node count} --- the largest single cost driver: the fast solve
        scales as $N^{1.32}$, the slow step linearly, and the boundary layer is
        125\,k of the \SI{1000}{\nano\meter} mesh's 130\,k tets. A modelling
        decision, not an optimization, but it moves runtime more than any code
        change here.
\end{enumerate}

\noindent
One knob is deliberately excluded: \texttt{fem\_every}. Splitting error is first
order in the fast-solve spacing, so time saved there comes straight out of the
physics.

\section*{Reproducing}
\begin{lstlisting}
# slow-step shard plan and the escape hatch
python -c "from dislocluster_code.zerod.cpp_bridge import _shard_plan; print(_shard_plan(190000))"
DISLOCLUSTER_SOLVER_SHARDS=1    # restores the single-process path exactly

# fast-step stage timers
wsl -e bash -c "OMP_NUM_THREADS=20 <build>/tools/DDomp/DDomp <simdir> 2>&1 | grep 'mobile stages'"
DISLOCLUSTER_DDOMP_THREADS=20   # overrides the launcher default
\end{lstlisting}

\paragraph{Caveat.} Every number here comes from one host. The processor-group
limit is Windows-specific, the shard plan is sized from
\texttt{os.cpu\_count()}, and the fast-step thread optimum already shifted
between the two problem sizes measured (40 at \SI{200}{\nano\meter}, 20 at
\SI{1000}{\nano\meter}). Single-run stage timings vary by up to $\sim$25\,\% on
this host, so the per-phase table of \S\ref{sec:retract} supports its
qualitative conclusion --- which phase scales --- but not its third digit.
Re-measure before sizing a campaign on different hardware.

\end{document}
